% ============================================================
% brand.tex — PLACEHOLDER brand identity (generic starting point)
% ============================================================
% Everything in THIS file is what makes a rendered book look like a
% specific brand: colors, fonts, and the book title used in the
% running header. preamble.tex (loaded after this file) is the
% structural engine — page layout, table mechanics, front/back
% matter switching, landscape handling — and doesn't change between
% brands or between books. This is the ONE file you need to edit to
% adapt this toolkit to a real company's brand; everything else in
% the project should work unmodified.
%
% Load order matters: this file must be listed BEFORE preamble.tex
% in _quarto.yml's `include-in-header:` list, since preamble.tex
% references the colors and \giBookTitle defined here.
%
% Every value below is a generic placeholder — swap them for the
% real brand's palette and typefaces. Search/replace the `gi` prefix
% on the color and macro names if you want to rename this toolkit's
% naming convention entirely (it's just a naming convention, not a
% requirement — the engine in preamble.tex references these exact
% names, so if you rename them, update preamble.tex's references too).
% ============================================================

\usepackage{fontspec}
\usepackage[table]{xcolor}

% ============================================================
% FONTS
% This starting point uses TeX Gyre Termes (body) + TeX Gyre Heros
% (headings) — professional, freely-licensed faces that ship with
% every standard TeX/TinyTeX install, so this toolkit renders with
% ZERO extra setup, on any machine, with no font files to source or
% bundle. Swap in a real brand's custom fonts when you have them:
%
%   1. Put the static-weight .ttf/.otf files in a fonts/ folder
%      (Regular, Bold, Italic, Bold Italic — variable-font files
%      don't work reliably with LaTeX's font engine).
%   2. Replace the two blocks below with this pattern:
%        \setmainfont{YourBodyFont}[Path=fonts/, Extension=.ttf, ...]
%        \newfontfamily\headingfont{YourHeadingFont}[Path=fonts/, ...]
%      Match the font FAMILY name you pass here to your files'
%      basename exactly (e.g. family "YourBodyFont" needs files
%      named YourBodyFont-Regular.ttf, not "Your Body Font-Regular")
%      — a mismatch here is a real, easy-to-hit error.
% ============================================================
\setmainfont{TeX Gyre Termes}

\newfontfamily\headingfont{TeX Gyre Heros}[Mapping = tex-text]
% Without Mapping=tex-text, chapter/section titles show a literal
% "---" instead of an em dash (—), and "--" instead of an en dash
% (–), even though body text renders them correctly. fontspec
% applies this dash/quote-conversion mapping automatically for
% \setmainfont (above) but NOT automatically for \newfontfamily-
% loaded fonts like this one — has to be requested explicitly here
% to match. Keep this option if you swap in a real brand's custom
% heading font later.

% ============================================================
% COLOR PALETTE — placeholder
% Define once, reuse everywhere — the equivalent of :root { --var } in CSS.
% Replace every hex value below with the real brand's palette.
% ============================================================
\definecolor{giPrimary}{HTML}{1E3A5F}      % primary heading / UI color
\definecolor{giPrimaryTint}{HTML}{E4EAF1}  % primary +90% — table header bg
\definecolor{giSecondary}{HTML}{4A4A68}    % secondary accent color
\definecolor{giCrimson}{HTML}{A6303F}      % core brand accent — titles, rules
\definecolor{giCrimsonTint}{HTML}{F7E5E7}  % accent +90% — subtle accent bg
\definecolor{giInk}{HTML}{1A1A1A}          % body text color
\definecolor{giWarning}{HTML}{E0954E}      % medium-risk / caution (data tables)
\definecolor{giWarningTint}{HTML}{FCF1E6}  % warning +90%
\definecolor{giNeutralWarm}{HTML}{B8B0A6}  % muted UI elements
\definecolor{giNeutralGray}{HTML}{DADCDC}  % borders/dividers
\definecolor{giBgTint}{HTML}{F1F3F7}       % near-white — alt row bg
\definecolor{giSuccess}{HTML}{3A9E82}      % low-risk / positive
\definecolor{giAccent}{HTML}{2A4FD6}       % links, highlights

% ---- Semantic alias: giCrimson doubles as the risk-table "danger" color ----
\colorlet{giDanger}{giCrimson}
\colorlet{giDangerTint}{giCrimsonTint}

% ---- Convenience aliases used elsewhere in this file ----
\colorlet{tableHeaderBg}{giPrimaryTint}
\colorlet{tableAltRow}{giBgTint}

% ---- Table shading (header / first column / alternating rows) ----
\definecolor{giTableHeaderShade}{HTML}{EAEAEC}   % table header row background
\definecolor{giTableFirstColShade}{HTML}{E9E9F0} % first-column background
\definecolor{giTableAltRowShade}{HTML}{F3F3F5}   % alternating body-row background
\definecolor{giTableDivider}{HTML}{DFE1E1}       % hairline row-divider color

% ============================================================
% BOOK TITLE — used in the running header (alternates with the
% current chapter name — see preamble.tex). This is the SERIES/
% PUBLICATION title, not the specific book's title — that lives in
% title-page-content.tex instead, since it'd change per book while
% this stays constant across a whole catalog/series.
% ============================================================
\newcommand{\giBookTitle}{Scaling Before Stability — Colophon Sample}
